\documentclass[12pt,a4paper]{article}
\usepackage[utf8]{inputenc}
\usepackage[spanish,es-noquoting]{babel}
\usepackage[T1]{fontenc}
\usepackage{lmodern}
\usepackage{geometry}
\geometry{top=2cm, bottom=2cm, left=2.5cm, right=2.5cm}
\usepackage{xcolor}
\usepackage{listings}
\usepackage{enumitem}
\usepackage{fancyhdr}
\usepackage{hyperref}
\usepackage{tikz}
\usetikzlibrary{arrows.meta, positioning, calc, shapes.geometric}

\definecolor{codegreen}{rgb}{0,0.6,0}
\definecolor{backcolour}{rgb}{0.95,0.95,0.92}

\lstset{
    language=C,
    backgroundcolor=\color{backcolour},
    commentstyle=\color{codegreen},
    keywordstyle=\color{blue},
    basicstyle=\ttfamily\small,
    breaklines=true,
    frame=single,
    numbers=left,
    tabsize=4,
    extendedchars=true,
    showstringspaces=false,
    literate={á}{{\'a}}1 {é}{{\'e}}1 {í}{{\'i}}1 {ó}{{\'o}}1 {ú}{{\'u}}1 {ñ}{{\~n}}1
}

\pagestyle{fancy}
\fancyhf{}
\lhead{Monitorías Sistemas Operativos 2026-I}
\rhead{Capítulo 2: Tuberías (Pipes)}
\cfoot{\thepage}

\title{\textbf{Taller de Lógica: Comunicación mediante Tuberías}}
\author{Malak Sanchez \\ \small Monitorías de Sistemas Operativos}
\date{\today}

\begin{document}

\maketitle

\section*{Introducción}
Este taller se enfoca en el uso de tuberías anónimas (\texttt{pipe}) para el flujo de datos entre procesos. Se explorarán conceptos como la redirección de descriptores de archivos (\texttt{dup2}), la ejecución de comandos externos (\texttt{exec}) y la comunicación bidireccional.

\section*{Bloque 1: Implementación de Tuberías Básicas}
\textit{Utilice fork() y pipe() para implementar flujos de datos entre procesos y redireccionar la salida estándar si es necesario.}

\subsection*{Ejercicio 1: Lectura Remota de Archivos}
Este ejercicio simula un servidor de archivos básico. El proceso padre actúa como cliente y el hijo como servidor.
\begin{enumerate}
    \item El Padre envía el nombre de un archivo (ej: \texttt{"datos.txt"}) al hijo a través de una tubería.
    \item El Hijo abre el archivo, lee su contenido y lo envía de vuelta al padre mediante una segunda tubería.
    \item El Padre imprime en pantalla el contenido recibido.
    \item \textbf{Reto:} El hijo debe enviar un mensaje de error si el archivo no existe.
\end{enumerate}

\begin{center}
\begin{tikzpicture}[node distance=1.5cm, box/.style={draw, rectangle, minimum width=2.5cm, minimum height=1cm, fill=blue!10}]
    \node[box] (P) {Padre (Cliente)};
    \node[box] (H) [right=of P, xshift=3cm] {Hijo (Servidor)};
    
    \draw[-Stealth, bend left] (P) to node[above] {Nombre del archivo} (H);
    \draw[-Stealth, bend left] (H) to node[below] {Contenido / Error} (P);
\end{tikzpicture}
\end{center}

\newpage

\section*{Bloque 2: Comunicación de Datos Dinámicos}
\textit{Desarrolle programas que intercambien información estructurada entre procesos.}

\subsection*{Ejercicio 2: El Eco Bidireccional}
Implemente una comunicación de "ida y vuelta" entre un padre y su hijo utilizando \textbf{dos tuberías}.
\begin{enumerate}
    \item El Padre lee una cadena por teclado y la envía por el \texttt{Pipe A}.
    \item El Hijo recibe la cadena, la convierte a MAYÚSCULAS y la envía por el \texttt{Pipe B}.
    \item El Padre recibe la respuesta y la imprime en pantalla.
\end{enumerate}

\begin{center}
\begin{tikzpicture}[node distance=2cm, proc/.style={draw, circle, minimum size=1.5cm, fill=orange!20}]
    \node[proc] (P) {Padre};
    \node[proc] (H) [right=of P, xshift=3cm] {Hijo};

    \draw[-Stealth, bend left] (P) to node[above] {Pipe A (Data)} (H);
    \draw[-Stealth, bend left] (H) to node[below] {Pipe B (Response)} (P);
\end{tikzpicture}
\end{center}

\newpage

\subsection*{Ejercicio 3: El Recolector de Estadísticas}
Un proceso padre crea $N$ procesos hijos. Cada hijo debe:
\begin{enumerate}
    \item Generar un número aleatorio entre 1 y 100.
    \item Enviar dicho número al padre a través de una tubería \textbf{compartida} o una tubería por cada hijo.
\end{enumerate}
El padre debe esperar a que todos los hijos terminen, recibir los números y calcular el promedio, el máximo y el mínimo.

\subsection*{Ejercicio 4: El Anillo de Procesos}
Cree una estructura de anillo con 4 procesos (A $\rightarrow$ B $\rightarrow$ C $\rightarrow$ D $\rightarrow$ A).
\begin{center}
\begin{tikzpicture}[node distance=2cm, proc/.style={draw, circle, minimum size=1cm, fill=green!10}]
    \node[proc] (A) {A};
    \node[proc] (B) [right=of A] {B};
    \node[proc] (C) [below=of B] {C};
    \node[proc] (D) [below=of A] {D};

    \draw[-Stealth] (A) -- (B); \draw[-Stealth] (B) -- (C);
    \draw[-Stealth] (C) -- (D); \draw[-Stealth] (D) -- (A);
\end{tikzpicture}
\end{center}
Un "token" numérico debe viajar por el anillo. Cada proceso debe sumarle 1 al token y pasarlo al siguiente. El juego termina cuando el token da 5 vueltas completas al anillo.

\section*{Bloque 4: Análisis de Código}
\textit{Analice el fragmento y responda: ¿Por qué el programa se queda "colgado" (deadlock) y no termina?}

\begin{lstlisting}
int fd[2];
pipe(fd);
if (fork() == 0) {
    char buf[10];
    // El hijo no cierra nada
    read(fd[0], buf, 10);
    printf("Hijo leyo: %s\n", buf);
    exit(0);
}
// El padre tampoco cierra fd[1] explicitamente antes de esperar
write(fd[1], "Hola", 5);
wait(NULL);
\end{lstlisting}

\end{document}
